\subsection{A finite upper saddle near ratio one}\label{rw:upper-resonance}

We first treat the finite hyperbolic saddle, then the repeated and weak
ends at infinity.  Near ratio one, expanding a power in the ratio
would lose uniformity when its logarithmic input is large.  Instead
we keep the exact compensator.  The original field is affine at zero
scale, so its nonlinear orbital remainder has a factor $\epsilon$,
which persists in the compensator and curvature estimates.  Fix a compact
interval $0<m_*\le\bar m\le m^*$, bounded remaining coefficients and a
small neighborhood on which its eigenvalue ratio satisfies
$|r-1|\le1/64$. The ratio used below is the actual ratio of the original
field. Dividing \eqref{rw:geometry:upper-field} by $Y+\epsilon>0$ gives
\begin{equation}\label{rw:upper:divided-field}
 \dot X=\frac{Y(\bar mY-1)+\epsilon(\bar BX^2+\bar aX)}{Y+\epsilon},
 \qquad
 \dot Y=X+\bar cY+\frac{\nu Y^2}{Y+\epsilon}.
\end{equation}
At $\epsilon=0$ this is an affine field. Its saddle is
$(-\bar d/\bar m,1/\bar m)$, where $\bar d=\bar c+\nu$, and its two
eigenvalues have opposite signs separated from zero. The implicit
function theorem gives the perturbed saddle and its eigendirections on
a common analytic coefficient neighborhood. Write its eigenvalues as
$\lambda_u>0>\lambda_s$ and set $r=-\lambda_s/\lambda_u$.

To bound the analytic jets near this finite saddle, choose a complex
neighborhood whose real part lies in $Y>Y_s/2>0$.  On a sufficiently
small such neighborhood, $Y+\epsilon$ stays away from zero.  In linear
eigenvector coordinates the stable and unstable graph equations have
coefficient denominators $n\lambda_u-\lambda_s$ and
$\lambda_u-n\lambda_s$, respectively. Their absolute values are at
least $c(n+1)$ for $n\ge1$ on a smaller complex coefficient disk.
The nonlinear part of \eqref{rw:upper:divided-field} is $\epsilon$
times an analytic function. On a sufficiently small state disk the
coefficient recursion is therefore a contraction in the norm
$\sum|a_n|\delta^n$; the factor $(n+1)^{-1}$ absorbs the derivative
index in the graph equation. This constructs both axes on one disk,
with their nonlinear parts divisible by $\epsilon$. Cauchy estimates
on smaller disks give any prescribed finite state and coefficient jets,
including Euler scale jets of the quotient by $\epsilon$.

Straighten these axes, orient their coordinates $U,V$ positively in the
quadrant of the local returning arc, and write the orbital equation as
$V_U=-(V/U)H(U,V)$, with $H(0,0)=r>0$. Separate positive changes on the
two axes, obtained by integrating
$d\log\widetilde U/d\log U=H(U,0)/r$ and
$d\log\widetilde V/d\log V=r/H(0,V)$, make the orbital coefficient
constant on each axis. Its difference from $r$ is then divisible by
$UV$. Dropping the tildes gives
\begin{equation}\label{rw:upper:axis-equation}
 \frac{dV}{dU}=-r\frac VU-V^2R(U,V),\qquad R=\epsilon\widetilde R,
\end{equation}
where $\widetilde R$ has common analytic bounds.  The local maps below
follow this positive quadrant.  To obtain the complete upper path,
they will be composed with its two regular connecting arcs.

The coordinates can be chosen to preserve the reversible component
$\bar a=\bar c=\nu=0$ exactly. Indeed, with $Y_s=1/\bar m$, put
\begin{gather*}
 W(Y)=\left(\frac{Y+\epsilon}{Y_s+\epsilon}\right)^{-2\epsilon\bar B},
 \qquad
 I(Y)=\int_{Y_s}^Y\frac{2s(\bar ms-1)}{s+\epsilon}W(s)\,ds.
\end{gather*}
Direct differentiation gives the first integral $WX^2-I$.
Writing $I=(Y-Y_s)^2J$ yields
$J(Y_s)=\bar m/(1+\epsilon\bar m)>0$. Thus
$\eta=(Y-Y_s)\sqrt{J/\bar m}$,
$\zeta=\sqrt{W/\bar m}X$ and
$U=-(\eta+\zeta)/2$, $V=-(\eta-\zeta)/2$ transform that integral to
$-4\bar mUV$. Extend these analytic coordinates to the coefficient
neighborhood before straightening the perturbed axes. The subsequent
axis changes are the identity on this component; there $r=1$ and
$R=0$ exactly.  This local first integral supplies the normalization
on the reversible component.

\begin{proposition}[The uniform resonant passage]\label{rw:upper:prepared-map}
After choosing common small local sections and positive coordinates on
them, the passage from $(t,1)$ to $(1,v)$ has the form
\begin{equation}\label{rw:upper:comp-map}
 v=t^r e^{w(t)},\qquad
 w=-\sigma t\omega(t;r-1)+R_w(t),\qquad
 \omega(t;\kappa)=\int_0^{\log(1/t)}e^{-\kappa s}\,ds.
\end{equation}
For each prescribed finite integer $N$, the neighborhoods can be chosen
so that, with $D=t\partial_t$,
\begin{equation}\label{rw:upper:remainder}
 |D^i\partial_{\mathbf b}^{\alpha}E_\epsilon^jR_w|
 \le C_N\epsilon t^{7/4},\qquad i+|\alpha|+j\le N,
 \qquad 0<t<e^{-6}.
\end{equation}
Ratio derivatives may also be included, counting towards the same total
order. The coefficients satisfy $\sigma=O(\epsilon)$ in these norms.
In particular $N=6$ supplies four input Euler derivatives and two
coefficient derivatives of $R_w$. The estimate places no bound on
$|(r-1)\log t|$.
\end{proposition}

\begin{proof}
First remove the nonconstant restrictions of $R$ to its axes without
dividing by $r-1$. Put $s_0=R(0,0)$ and, writing
$R(U,0)=s_0+\sum_{n\ge1}R_{n0}U^n$, set
\[
 f(U)=-\sum_{n\ge1}\frac{R_{n0}U^{n+1}}{n+1-r},\qquad
 V_1=\frac V{1+f(U)V}.
\]
A direct differentiation gives
\[
 R_1(U,V_1)=R\left(U,\frac{V_1}{1-f(U)V_1}\right)
                  +f'(U)-rf(U)/U,
 \qquad R_1(U,0)=s_0.
\]
Next set $U_1=U/(1+g(V_1)U)$, where $g(0)=g'(0)=0$ solves
\[
 g'(V)-\frac{g(V)}{rV}=\frac{R_1(0,V)-s_0}{r^2}.
\]
The new axis restriction is
$R_2(0,V)=R_1(0,V)-r^2g'(V)+rg(V)/V=s_0$.
Since $g=O(V^2)$, the previously normalized axis is unchanged.
The denominators of the two series are $n+1-r$ and $n+1-1/r$,
$n\ge1$. On $|r-1|\le1/64$ they have a fixed positive margin and
grow linearly with $n$. The changes and their inverses therefore have
uniform analytic bounds on smaller disks, and $f,g$ retain their factor
$\epsilon$. Exact division of $R_2-s_0$ by $U_1V_1$, followed by one
common positive rescaling of both sections, gives
\begin{equation}\label{rw:upper:resonant-orbit}
 V_U=-rV/U-\sigma V^2-UV^3\mathcal R(U,V),\qquad
 |\sigma|+\|\mathcal R\|\le C\epsilon.
\end{equation}
The bound includes whatever finite jets are reserved before this
construction. The two shears are also the identity on the reversible
component considered above.

The reciprocal equation before the shears has the exact integral
representation
\[
 V(U)=\frac{(t/U)^r}{1+t^r\int_t^U s^{-r}R(s,V(s))\,ds}.
\]
It proves positivity and existence on the full rectangle for small
$\epsilon$. In \eqref{rw:upper:resonant-orbit} put
$U=te^\theta$, $V=e^{-r\theta}\zeta$, and
$L=\log(1/t)$. Then, exactly,
\begin{equation}\label{rw:upper:zeta}
 \zeta_\theta=-\sigma te^{-(r-1)\theta}\zeta^2
 -t^2e^{-2(r-1)\theta}\zeta^3
        \mathcal R(te^\theta,e^{-r\theta}\zeta),\qquad \zeta(0)=1.
\end{equation}
With $\delta=1/64$ the two kernel integrals are bounded respectively by
$t^{1-\delta}L$ and $t^{2-2\delta}L$. A bootstrap on the integral
equation gives $1/2<\zeta<2$ and
$\zeta-1=O(\epsilon t^{1-\delta}(1+L))$ on $[0,L]$.
Integrating its logarithmic derivative now yields
\begin{align*}
 R_w={}&-\sigma t\int_0^L e^{-(r-1)\theta}(\zeta-1)\,d\theta\\
 &-t^2\int_0^L e^{-2(r-1)\theta}\zeta^2
                  \mathcal R(te^\theta,e^{-r\theta}\zeta)\,d\theta.
\end{align*}
Consequently $|R_w|\le C\epsilon t^{2-2\delta}(1+L)^2$.

Differentiate \eqref{rw:upper:zeta} at fixed $\theta$. Ratio derivatives
add polynomial factors in $\theta$; input Euler derivatives preserve
$te^\theta$. The highest variation has the same small-$L^1$ linear
kernel, so induction bounds it by the displayed powers of $t$ times a
fixed power of $1+L$. Endpoint rows use $DL=-1$ and retain those powers.
In the first integral for $R_w$, all variations of $\zeta-1$ retain
their first-order smallness. Therefore
\[
 |D^i\partial_{\mathbf b}^{\alpha}E_\epsilon^jR_w|
 \le C_N\epsilon t^{2-2/64}(1+L)^{M_N}.
\]
The strict exponent margin $2-2/64-7/4=7/32$ absorbs every such fixed
logarithmic power and proves \eqref{rw:upper:remainder}. Ordinary
$\epsilon$ derivatives can differentiate away the front factor
$\epsilon$ and are not covered by that displayed conclusion.
\end{proof}

The curvature formula follows from the exact identity
\[
 E:=t\frac{v''}{v'}-(r-1)=Dw+\frac{D^2w}{r+Dw}.
\]
Inserting \eqref{rw:upper:comp-map} gives
\begin{equation}\label{rw:upper:curvature}
 E=-\sigma Q+R_E,\qquad
 Q=2rt\omega(t;r-1)-(2+1/r)t,\qquad Q''=-2r^2t^{r-2}.
\end{equation}
Here $|D^i\partial_{\mathbf b}^{\alpha}E_\epsilon^jR_E|
\le C\epsilon t^{7/4}$ for $i+|\alpha|+j\le N-2$ if $N$ orders were
retained in \eqref{rw:upper:remainder}. Indeed the error in replacing
$(r+Dw)^{-1}$ by $1/r$ is a product of two first corrections, bounded
by the same $t^{2-2/64}(1+L)^M$ estimate. Thus a curvature use requiring
four Euler and two coefficient derivatives reserves $N=8$, whereas
four Euler and two coefficient derivatives of $w$ itself require
$N=6$. Since $\omega\ge\int_0^6e^{-s/64}\,ds>5$ and
$Q'=2r^2\omega-2r-2-1/r$, both $Q$ and $Q'$ are positive on
$0<t<e^{-6}$ for the stated ratio interval.  Both estimates retain
$t^{r-1}$ exactly and therefore remain valid when
$|(r-1)\log t|$ is large.

The finite connecting arcs contribute a separate regular factor. On a
small low-height segment where $|X|\ge c>0$, the denominator of
$dX/dY$ obtained from \eqref{rw:upper:divided-field} is bounded below
in modulus by $c(Y+\epsilon)$. The numerator and all its fixed $X$
and coefficient derivatives are $O(Y+\epsilon)$; Euler scale
differentiation preserves this property. The equation in $Y$ and its
finite variational equations consequently have bounded coefficients.
Comparing it with its zero-scale equation leaves only the integrable
layer carrier $\epsilon/(Y+\epsilon)$, whose integral is
$\epsilon\log(1+\delta/\epsilon)$. Beyond this segment the vector
field is regular on a compact arc. These arguments give the required
finite section and coefficient jets, with zero-scale error
$O(\epsilon(1+|\log\epsilon|)+|\nu|)$ when the comparison also sets
$\nu=0$.  The scale derivatives in this estimate are Euler
derivatives; ordinary $\epsilon$ derivatives of the error may be
unbounded.

\subsection{Repeated and weak upper passages}\label{rw:upper-degenerate-preparation}

As an upper root becomes repeated or weak, the preceding hyperbolic
preparation ceases to be uniform.  We instead use a small angular
width to make the slow variational growth weaker than the normal
contraction.  In projective coordinates
$R=1/Y$, $\varphi=X/Y$, put $b=1-\epsilon\bar B$ and
$\bar d=\bar c+\nu$. A positive projective time change gives the exact
field
\begin{equation}\label{rw:upper:projective}
\begin{split}
 R'&=-R(\varphi+\bar d)-\epsilon R^2(\varphi+\bar c),\\
 \varphi'&=\bar m-b\varphi^2-\bar d\varphi
       +R[-1+\epsilon\varphi(\bar a-\bar c-\varphi)].
\end{split}
\end{equation}
We describe separately the quadratic angular minimum and the weak
radial factor. In both constructions all finite derivative orders are
fixed before shrinking the strip.

\paragraph{The quadratic angular factor.}
Fix a limiting repeated field $\bar c=2a_*>0$,
$\bar m=-a_*^2$, $\epsilon=\nu=0$, with bounded $\bar B,\bar a$.
The symbol $a_*$ here is a positive constant, unrelated to $\bar a$.
On $R=0$ set
\[
 v_0=-\varphi-\frac{\bar d}{2b},\qquad
 \mu_U=-\frac{\bar d^2+4b\bar m}{4b^2},\qquad
 \Xi=\frac{\bar d(1-2\epsilon\bar B)}{2b}.
\]
Dividing the vector field by the positive factor $b$ makes its angular
restriction exactly $v_0'=\mu_U+v_0^2$, and its normal rate is
$(\Xi-v_0)/b$. Near the chosen field $\Xi/b$ has a fixed positive
margin. The line $R=0$ is exactly invariant.

Choose an angular outer width $\varpi>0$. Extend $v_0^2$ outside
$|v_0|<\varpi$ to a smooth bounded function with first derivative
$O(\varpi)$, and extend the normal coefficients preserving positive
margins for $\beta$ and $\beta+R\beta_R$. The extension is made on a
larger real rectangle than the one used for the original passage.
Since its scalar derivative is $O(\varpi)$, the flow induction
\eqref{rw:transitions:slow-variations}, now with $q=1$, gives
$C(1+t)^Me^{C_N\varpi|t|}$ through order $N+2$, also for the inverse
flow and parameter derivatives of $\mu_U$. Reserve all products of
these and the exit-time jets in $C_N$, then choose
$C_N\varpi<\beta_*/16$, followed by radial width and coefficient
neighborhood. The angular equation has no small factor $\epsilon$.

Let $\phi_t$ denote the complete scalar flow on $R=0$ and let
$(v_0(t),R(t))$ be a trajectory of the complete auxiliary field. The
phase integral
\[
 v=\lim_{t\longrightarrow\infty}\phi_{-t}(v_0(t))
\]
exists by the inverse-flow construction of Proposition~\ref{rw:phase}:
the forcing $RO(1)$ gives the majorant
$C(1+t)^Me^{-(\beta_*-C_N\varpi)t}$. The limit equals $v_0$ on the
axis and obeys its scalar flow identity; hence $v=v_0+RO_{C^N}(1)$
is a coordinate. Reusing $R$ for a positive normal coordinate, the original
field, on the smaller rectangle, is exactly
\begin{equation}\label{rw:upper:repeated-normal}
 v'=\mu_U+v^2,
 \qquad R'=-R\beta_U(v,R),\qquad
 \beta_U(v,0)=\frac{\Xi-v}{b},\qquad
 \beta_U,\ \beta_U+R\beta_{U,R}>\beta_*>0.
\end{equation}
Shrinking the radial rectangle retains the normal margins. Together
with the angular variation bound and the fixed transverse exit below,
these estimates give the hypotheses of
Proposition~\ref{rw:relative-core} with $q=1$. The small angular width
provides the required domination in this upper passage.

For a fixed small radial input $R_0>0$ and a fixed positive angular
exit $v_R<\varpi/2$, restrict $|\mu_U|<v_R^2/2$. The exit speed
is then at least $v_R^2/2$. The actual outgoing normal value is
\begin{equation}\label{rw:upper:repeated-output}
 G(v)=R_0\exp\left[-\int_v^{v_R}
       \frac{\beta_U(s,0)}{\mu_U+s^2}\,ds\right]e^{j_U(v)}.
\end{equation}
Its relative amplitude $j_U$ has the finite mixed bounds of
Proposition~\ref{rw:relative-core} and
Lemma~\ref{rw:transitions:moving-rate}, on the stated auxiliary
domain. The physical domain is the whole-path condition
\begin{equation}\label{rw:upper:quadratic-passing}
 \mu_U+s^2>0\quad(v\le s\le v_R),
 \quad\text{equivalently}\quad
 \mu_U>0\ \text{or}\
 [\mu_U\le0\ \text{and}\ v>\sqrt{-\mu_U}],
 \qquad v<v_R.
\end{equation}
A scalar root cannot be crossed in finite time.  The exact phase
coordinate carries this obstruction to every stable fibre of the
original field.  A fixed passage from
$-v_R$ to $v_R$ therefore requires $\mu_U>0$.  The full-interval
condition is needed even when the slow factor is positive at both
endpoints.

To apply this preparation to the upper comparison, we next locate the
entry reached from the finite section, on both sides of its limiting
threshold.  At the limiting repeated field the affine upper
system in \eqref{rw:upper:divided-field} has the exact coordinates
\[
 V=Y+a_*^{-2},\qquad Z=X+a_*Y-a_*^{-1},\qquad r_0=Z/V,
\]
which satisfy $V'=a_*V+Z$, $Z'=a_*Z$, and $r_0'=-r_0^2$.
For an orbit starting at $(p,0)$, put $P=1/a_*$ and $Z_0=p-P$.
At a fixed large $V=V_e$ reached at time $T$,
\[
 V_e=e^{a_*T}(a_*^{-2}+Z_0T),\qquad
 \frac{\partial r_{0,e}}{\partial p}
 =\frac{e^{a_*T}}{V_e}
    \frac{p}{a_*Z_0T+Z_0+a_*^{-1}}>0.
\]
In particular the derivative at $p=P$ is $a_*^2$.
Choose $v=-r_0$ and normal coordinate $1/V$ at this limiting field.
Then \eqref{rw:upper:repeated-normal} is exact with
$\beta_U=a_*-v$. Choose the auxiliary extension in these exact
limiting coordinates, and extend this choice smoothly in the parameters;
thus $j_U=0$ there. On each fixed compact finite entry
interval about $P$ on which the connecting orbit reaches this section
transversally, persistence and the low-layer estimate above give
$v_e'(p)<-c<0$ and bounded finite section and coefficient jets.
The limiting entry is zero at $P$ and negative for $p>P$; such negative
entries are retained and pass precisely when
\eqref{rw:upper:quadratic-passing} permits them. By the limiting
normalization and continuity, $j_U$ tends to zero in the reserved norms
as the coefficients tend to this limiting field.

\paragraph{The weak radial factor.}
Fix $\bar c_0>0$, $\bar m=0$, and $\epsilon=\nu=0$.
On a sufficiently small coefficient neighborhood define the actual
quantities
\begin{equation}\label{rw:upper:weak-parameters}
\begin{split}
 \Delta&=\sqrt{\bar d^2+4b\bar m},\qquad
 \varphi_s=\frac{-\bar d-\Delta}{2b},\\
 \gamma&=\varphi_s+\bar d
 =-\frac{2(\bar m+\epsilon\bar B\bar d^2)}
              {\bar d(1-2\epsilon\bar B)+\Delta}.
\end{split}
\end{equation}
The denominator is positive and bounded away from zero. In particular
the finite-scale weak face is
$\bar m+\epsilon\bar B\bar d^2=0$; replacing it by $\bar m=0$
would move the radial root. Reverse the projective time in
\eqref{rw:upper:projective} and put $y=\varphi-\varphi_s$.
The exact system is
\begin{equation}\label{rw:upper:weak-original}
\begin{split}
 \dot R&=R(\gamma+y)+\epsilon R^2(\varphi_s+y+\bar c),\\
 \dot y&=-\Delta y+by^2
 +R\{1-\epsilon(\varphi_s+y)(\bar a-\bar c-\varphi_s-y)\}.
\end{split}
\end{equation}
Its normal eigenvalue has a fixed negative margin and its slow
radial derivative is small when $|\gamma|$, the radial width and the
normal width are small.

We construct the weak graph with the prescribed finite smoothness on
a common real domain. Extend \eqref{rw:upper:weak-original} smoothly
outside a slightly larger rectangle, preserving $R=0$, normal
contraction and the small radial derivative. The time-one graph
transform has base variation at most $e^{C\varpi}$ and normal
variation at most $e^{-\Delta_*/2}$ after reducing the rectangle.
Choose its width so that
$e^{-\Delta_*/2}e^{(N+2)C\varpi}<1$ and the bound
$C_N\varpi<\beta_*/16$ holds for all the finite products reserved
above.
The mixed graph-transform induction of Proposition~\ref{rw:phase}, with
this contracting highest-jet multiplier, constructs $y=h(R)$,
$h(0)=0$, with the reserved mixed bounds on the common auxiliary domain.

The linear graph equation gives
\[
 h'(0)=\alpha=
 \frac{1-\epsilon\varphi_s(\bar a-\bar c-\varphi_s)}{\Delta+\gamma}.
\]
The reduced equation is $\dot R=RG(R)$, where
\begin{equation}\label{rw:upper:weak-G}
 G(R)=\gamma+h(R)+\epsilon R(\varphi_s+h(R)+\bar c),\qquad
 G'(0)=\alpha+\epsilon(\varphi_s+\bar c)\longrightarrow1/\bar c_0>0.
\end{equation}
The same inverse-flow phase integral just used for the repeated strip,
now relative to this reduced equation, gives a coordinate
$\widetilde R$ constant on its strong fibres and equal to $R$ on the
weak graph. Its normal forcing is divisible by $R$, since $R=0$ is
invariant in \eqref{rw:upper:weak-original}. Thus
$\widetilde R=R$ times a positive unit on the smaller rectangle.
In particular physical $R>0$ corresponds exactly to
$\widetilde R>0$. The mixed convergence follows from the normal
margin and the chosen small slow width, as above.

Define
\[
 u=G(\widetilde R)-\gamma=\widetilde R\widetilde G(\widetilde R),
 \qquad \widetilde G(R)=\int_0^1G'(sR)\,ds>c>0.
\]
The positive time factor $G'/\widetilde G$ has bounded finite jets and
is equal to one at $R=0$. Dividing the vector field by this factor,
and straightening the weak graph, gives exactly
\begin{equation}\label{rw:upper:weak-normal}
 \dot u=u(\gamma+u),\qquad \dot Z=-Z\beta(u,Z),\qquad
 \beta(0,0)=\Delta,\qquad
 \beta,\ \beta+Z\beta_Z>\beta_*>0.
\end{equation}
Only one derivative of $G$ was spent here; it was reserved before the
graph construction.  The change is regular across $\gamma=0$.
Its slow derivative $\gamma+2u$ is $O(\varpi)$ on the extended small
strip, while its normal rate has the stated positive margin. All
hypotheses of Proposition~\ref{rw:relative-core} are therefore again
verified for $q=1$.

Fix a positive weak exit $u_R$ and count inputs below
$u_0<u_R/2$, taking $|\gamma|<u_R/2$. The actual passage domain is
\begin{equation}\label{rw:upper:weak-passing}
 \max(0,-\gamma)<u<u_0.
\end{equation}
This is precisely positivity of $u(\gamma+u)$ on the entire path to
$u_R$. At $\gamma<0$ its positive root $-\gamma$ blocks every entry
on its left. The phase projection and normal contraction give the
same obstruction for the original two-dimensional field.

The finite outgoing normal value, followed by its regular connector
to $Y=0$, is the inverse upper map. Choose its positive displacement
$t$ from the weak intersection $\hat p$ so that $p=\hat p-t(u)$.
The limiting weak graph is $X+\bar c_0Y=1/\bar c_0$; hence
$\hat p\to1/\bar c_0$, and its finite continuation is transverse.
The low-layer equation above bounds all the required section and
coefficient jets of this continuation. With $f=u(\gamma+u)$, the
relative normal formula and this connector give
\begin{equation}\label{rw:upper:weak-rate}
 t,t'>0,\qquad
 \frac{t'}t=\frac{N_U}{f},\qquad
 \frac{t''}{t'}=\frac{A_U}{f},\qquad
 N_U=\beta(u,0)+f\widetilde j'(u),\qquad
 A_U=N_U-f'+f\frac{N_U'}{N_U}.
\end{equation}
Here $\widetilde j$ includes the logarithm of the positive regular
connector unit and has the finite bounds from
Lemma~\ref{rw:transitions:moving-rate}. Since $f$ and $f'$ are
uniformly small, $N_U,A_U>c>0$ on a sufficiently small input and
coefficient neighborhood. The absolute output and all its selected
finite derivatives tend to zero uniformly as these neighborhoods
shrink: the travel time diverges, and normal decay dominates the
finite slow variational growth.

\paragraph{Moving inputs and lost exits.}
The phase maps were constructed on larger domains.  Their bounded
mixed jets therefore also control their restrictions to moving
transverse sections by composition.  For a full
repeated strip,
$T=\int_{-v_R}^{v_R}(\mu_U+s^2)^{-1}\,ds$ and
$|\partial_{\mu_U}^jT|\le C_j(1+T)^{2j+1}$.
For moving entries and the weak factor, apply the fixed-exit calculation
of Appendix~\ref{rw:app:relative-variations} with $q=1$, the stated
positive exit speeds and slow exponent $C_N\varpi$. This verifies
Proposition~\ref{rw:relative-core}; its extensions and
Corollary~\ref{rw:transitions:relative-continuity} apply in the reserved
mixed orders. The physical domains
remain \eqref{rw:upper:quadratic-passing} and
\eqref{rw:upper:weak-passing}, with the stated regular connections.
